
\section*{13. Flat affine groups over a regular base of
  dimension at most \(2\)}
\phantomsection\label{sec:sga3-VIB-13}%
\addcontentsline{toc}{section}{13. Flat affine groups over a
  regular base of dimension at most 2}

Let us begin by observing\footnote{\label{note:sga3-VIB-139}%
Editor's note: this section has been added.} that the familiar
argument showing that every affine algebraic group over a field
\(k\) is linear, as well as
\hyperref[lem:sga3-VIB-11.12]{Lemma 11.12}, extends to a group
scheme \(G\) that is affine, flat, and of finite type over a
Noetherian regular base scheme \(S\) of dimension at most \(2\).
For \(S\) of dimension at most \(1\), this proves
\hyperref[concl:sga3-VIB-11.11.1-b]{Remark 11.11.1(b)}. The
extension to \(\dim S=2\) rests on the following lemma,
communicated to us by O.~Gabber.

\medskip
\phantomsection\label{lem:sga3-VIB-13.1}%
\noindent\textbf{Lemma 13.1.---}
\emph{Let \(S\) be a normal Noetherian scheme, let
\(\mathcal A\) be a flat quasi-coherent
\(\mathcal O_S\)-module, let \(\mathcal F\) be a finite-type
\(\mathcal O_S\)-submodule of \(\mathcal A\), let
\(\mathcal F^{**}\) be its bidual, and let \(U\) be the
flatness locus of \(\mathcal F\), that is, the set of points
\(s\in S\) such that \(\mathcal F_s\) is a flat
\(\mathcal O_{S,s}\)-module.}
\begin{enumerate}
\item[\textup{(i)}]\phantomsection\label{concl:sga3-VIB-13.1-i}
  \emph{\(U\) is open in \(S\), and}
  \[
    \mathcal F^{**}=j_*j^*(\mathcal F),
  \]
  \emph{where \(j:U\hookrightarrow S\) denotes the
  inclusion.}%
  \begingroup
  \renewcommand{\thefootnote}{\fnsymbol{footnote}}%
  \footnote[1]{\label{note:sga3-VIB-13.1-i-U-ambient}
  Translator's note: the French re-edition prints
  \(j:U\hookrightarrow X\) in the statement. The ambient
  scheme throughout the lemma is \(S\), and its proof explicitly
  defines \(j:U\hookrightarrow S\); the latter is used here.
  Reinhold's English translation keeps the printed \(X\).}%
  \endgroup

\item[\textup{(ii)}]\phantomsection\label{concl:sga3-VIB-13.1-ii}
  \emph{For every quasi-coherent \(\mathcal O_U\)-module
  \(\mathcal E\), the canonical morphism}
  \phantomsection\label{form:sga3-VIB-13.1-projection}%
  \[
    j_*(\mathcal E)\otimes_{\mathcal O_S}\mathcal A
      \longrightarrow
    j_*\bigl(\mathcal E\otimes_{\mathcal O_U}j^*(\mathcal A)\bigr)
  \]
  \emph{is an isomorphism.}

\item[\textup{(iii)}]\phantomsection\label{concl:sga3-VIB-13.1-iii}
  \emph{In particular,}
  \[
    \mathcal V=j_*j^*(\mathcal F)
      \subset \mathcal A=j_*j^*(\mathcal A),
  \]
  \emph{and the canonical morphism}
  \phantomsection\label{form:sga3-VIB-13.1-tensor}%
  \[
    \mathcal V\otimes_{\mathcal O_S}\mathcal A
      \longrightarrow
    j_*j^*\bigl(\mathcal F\otimes_{\mathcal O_S}\mathcal A\bigr)
  \]
  \emph{is an isomorphism.}
\end{enumerate}

\emph{Proof.}
Replacing \(S\) by one of its connected components, we may
suppose that \(S\) is integral. By EGA~IV\(_2\), 2.1.12, the
flatness locus of \(\mathcal F\), namely the set of points
\(s\in S\) such that \(\mathcal F_s\) is a flat
\(\mathcal O_{S,s}\)-module, is an open subset \(U\) of \(S\);
write \(j:U\hookrightarrow S\) for its inclusion. Since
\(\mathcal A_s\) is flat, hence torsion-free, the same is true
of \(\mathcal F_s\). Thus \(U\) contains every point of
codimension at most \(1\). Consequently,
by \hyperref[bib:sga3-VIB-BAC]{\textup{[BAC]}}, VII, \S~4.2,
Corollary to Theorem~1,
\[
  \mathcal F^{**}=j_*j^*(\mathcal F),
\]
and hence there is a monomorphism
\[
  \mathcal F^{**}\longrightarrow j_*j^*(\mathcal A).
\]

The proof of \hyperref[concl:sga3-VIB-13.1-ii]{\textup{(ii)}}
is analogous to that of EGA~III, 1.4.15, recalled in
\hyperref[par:sga3-VIB-11.0]{11.0}. Since \(S\) is normal,
the morphism
\[
  \mathcal O_S\longrightarrow j_*j^*(\mathcal O_S)
\]
is an isomorphism \textup{(cf.\ EGA~IV\(_2\), 5.8.6 and
5.10.5)}. Applying
\hyperref[concl:sga3-VIB-13.1-ii]{\textup{(ii)}} to
\(\mathcal E=\mathcal O_U\) therefore gives
\[
  \mathcal A=j_*j^*(\mathcal A).
\]
Finally, since
\[
  j^*(\mathcal F\otimes_{\mathcal O_S}\mathcal A)
  =
  j^*(\mathcal F)\otimes_{\mathcal O_U}j^*(\mathcal A),
\]
the last assertion of
\hyperref[concl:sga3-VIB-13.1-iii]{\textup{(iii)}} follows
from \hyperref[concl:sga3-VIB-13.1-ii]{\textup{(ii)}} applied
to \(\mathcal E=j^*(\mathcal F)\). This proves the lemma.

Recall also that a finitely generated \(R\)-module \(M\) is
called \emph{reflexive} if the canonical morphism from \(M\)
to its bidual \(M^{**}\) is an isomorphism. If \(R\) is a
Noetherian regular ring of dimension at most \(2\), this
implies that \(M\) is projective. Indeed, for every finitely
generated \(R\)-module \(N\), choose a resolution
\[
  L_1\longrightarrow L_0\longrightarrow N\longrightarrow0,
\]
where \(L_0\) and \(L_1\) are finite free \(R\)-modules.
There is then an exact sequence
\phantomsection\label{seq:sga3-VIB-13.1-dual-resolution}%
\[
  0\longrightarrow N^*
    \longrightarrow L_0^*
    \longrightarrow L_1^*
    \longrightarrow Q
    \longrightarrow0,
\]
where \(Q\) is the cokernel of \(L_0^*\to L_1^*\). Since
\(R\) has homological dimension at most \(2\)
\textup{(\hyperref[bib:sga3-VIB-BAC]{[BAC]}, X, \S~\SGARefLink{prop:sga3-VIB-4.2}{4.2},
Corollary~1 to Theorem~1)}, it follows that \(N^*\) is
projective.

\medskip
\phantomsection\label{prop:sga3-VIB-13.2}%
\noindent\textbf{Proposition 13.2.---}
\emph{Let \(S\) be a Noetherian regular scheme of dimension at
most \(2\), let \(G\) be an affine flat \(S\)-group, and let
\(\mathcal A(G)\) be its affine algebra.}
\begin{enumerate}
\item[\textup{(i)}]\phantomsection\label{concl:sga3-VIB-13.2-i}
  \emph{If \(G\) is of finite type over \(S\), it is
  isomorphic to a closed subgroup of}
  \[
    H=\operatorname{Aut}_{\mathcal O_S}(\mathcal V)
  \]
  \emph{for some locally free \(\mathcal O_S\)-module
  \(\mathcal V\) of finite rank. If, moreover, \(S\) is
  affine, one may take
  \(\mathcal V=\mathcal O_S^{\oplus d}\) for some \(d\), so
  that \(H=\operatorname{GL}_{d,S}\).}

\item[\textup{(ii)}]\phantomsection\label{concl:sga3-VIB-13.2-ii}
  \emph{\(\mathcal A(G)\) is the filtered direct limit of
  flat finite-type \(\mathcal O_S\)-sub-Hopf algebras.}
\end{enumerate}

\emph{Proof.}
Let \(\mathcal B\) be a finite-type
\(\mathcal O_S\)-subalgebra of \(\mathcal A(G)\). Since every
coherent module on an open subset of \(S\) extends to a
coherent module on \(S\) \textup{(EGA~I, 9.4.5)}, there is a
coherent \(\mathcal O_S\)-submodule \(\mathcal M\) of
\(\mathcal B\) that generates \(\mathcal B\) as an
\(\mathcal O_S\)-algebra \textup{(\emph{loc.\ cit.}, 9.6.5)}.
By \hyperref[prop:sga3-VIB-11.10.bis]{11.10.bis},
\(\mathcal M\) is contained in a coherent \(G\)-stable
\(\mathcal O_S\)-submodule \(\mathcal F\).
\textup{(N.B. Since \(G\) is affine over \(S\), the proof of
\emph{loc.\ cit.} is simpler here: one may replace
\(f_*f^*(\mathcal E)\) by
\(\mathcal E\otimes_{\mathcal O_S}\mathcal A(G)\), and so
on.)}

Let \(j:U\hookrightarrow S\) be the inclusion, where \(U\)
is the flatness locus of \(\mathcal F\). By
\hyperref[lem:sga3-VIB-13.1]{Lemma 13.1} and the preceding
observation,
\[
  \mathcal V=j_*j^*(\mathcal F)
\]
is a locally free \(\mathcal O_S\)-submodule of
\(\mathcal A(G)\), and the canonical morphism
\phantomsection\label{form:sga3-VIB-13.2-V-stability}%
\[
  \mathcal V\otimes_{\mathcal O_S}\mathcal A(G)
    \longrightarrow
  j_*j^*\bigl(\mathcal F\otimes_{\mathcal O_S}\mathcal A(G)\bigr)
\]
is an isomorphism. Thus \(\mathcal V\) is a \(G\)-submodule
of \(\mathcal A(G)\). The action of \(G\) on \(\mathcal V\)
induces a morphism of affine \(S\)-groups
\[
  \rho_{\mathcal V}:
  G\longrightarrow
  H=\operatorname{Aut}_{\mathcal O_S}(\mathcal V)
\]
and therefore a morphism of \(\mathcal O_S\)-Hopf algebras
\[
  \varphi_{\mathcal V}:\mathcal A(H)\longrightarrow\mathcal A(G).
\]
Write \(\mathcal A_{\mathcal V}\) for the image of
\(\varphi_{\mathcal V}\). It is the affine algebra of a
closed subgroup \(G_{\mathcal V}\) of \(H\), the closed image
of \(\rho_{\mathcal V}\). We shall show that
\(\mathcal A_{\mathcal V}\) contains \(\mathcal B\).

The question is local on \(S\), so we may suppose that
\(S=\operatorname{Spec}(R)\) and that
\[
  V=\Gamma(S,\mathcal V)
\]
is a free \(R\)-module with basis \(v_1,\ldots,v_n\).
Then \(H\simeq\operatorname{GL}_{n,R}\), and
\(\mathcal A(H)\) is generated as an \(R\)-algebra by the
``matrix coefficients'' \(c_{ij}\) and by \(d^{-1}\), where
\(d\) denotes the determinant. Let \(\Delta\)
\textup{(respectively, \(\varepsilon\))} denote the
comultiplication \textup{(respectively, augmentation)} of
\(A=\Gamma(S,\mathcal A(G))\). For \(j=1,\ldots,n\), write
\phantomsection\label{form:sga3-VIB-13.2-matrix-coefficients}%
\[
  \Delta(v_j)=\sum_{i=1}^{n}v_i\otimes a_{ij}.
\]
Then \(a_{ij}=\varphi(c_{ij})\) belongs to
\(\operatorname{Im}(\varphi)\). Since \(V\) is an
\(R\)-submodule of \(A\), the identity
\[
  (\varepsilon\otimes\mathrm{id}_A)\circ\Delta
    =\mathrm{id}_A
\]
gives
\[
  v_j=\sum_i\varepsilon(v_i)a_{ij}
    \in\operatorname{Im}(\varphi).
\]
As \(V\) contains a set of generators of
\(B=\Gamma(S,\mathcal B)\), it follows that
\[
  B\subset\operatorname{Im}(\varphi).
\]

If \(G\) is of finite type over \(S\), take
\(\mathcal B=\mathcal A(G)\). Then \(\varphi\) is
surjective, so the morphism
\[
  G\longrightarrow
  H=\operatorname{Aut}_{\mathcal O_S}(\mathcal V)
\]
is a closed immersion.

If \(S\) is also affine, there is a locally free
\(\mathcal O_S\)-module \(\mathcal V'\) of finite rank such
that
\[
  \mathcal V\oplus\mathcal V'
    =\mathcal O_S^{\oplus d}.
\]
Regard \(\mathcal V'\) as a trivial \(G\)-module and replace
\(\mathcal V\) by \(\mathcal O_S^{\oplus d}\). This realizes
\(G\) as a closed subgroup of \(\operatorname{GL}_{d,S}\).

Finally, return to an arbitrary affine flat \(S\)-group
\(G\). By EGA~I, 9.4.9, \(\mathcal A(G)\) is the union of
its coherent \(\mathcal O_S\)-submodules \(\mathcal M\),
hence also of the \(\mathcal O_S\)-sub-Hopf algebras
\(\mathcal A_{\mathcal V}\) constructed above. This proves
\hyperref[concl:sga3-VIB-13.2-ii]{\textup{(ii)}}.

\medskip
\phantomsection\label{ex:sga3-VIB-13.3}%
\noindent\textbf{Example 13.3.---}
\emph{Let \(R\) be a discrete valuation ring with uniformizer
\(\pi\) and fraction field \(K\). Consider the filtered
projective system of \(R\)-groups}
\begin{center}
  \phantomsection\label{diag:sga3-VIB-13.3-projective-system}%
  \begin{tikzcd}[column sep=large]
    \cdots \arrow[r] &
    \mathbb G_{a,R} \arrow[r,"\times\pi"] &
    \mathbb G_{a,R} \arrow[r,"\times\pi"] &
    \mathbb G_{a,R}
  \end{tikzcd}
\end{center}
\emph{corresponding to the direct system}
\[
  R[X_0]\longrightarrow R[X_1]\longrightarrow R[X_2]
    \longrightarrow\cdots,
\]
\emph{whose transition morphisms are given by
\(X_n=\pi X_{n+1}\). Its projective limit \(G\) is a flat
affine \(R\)-group scheme, not of finite type, with trivial
special fiber and generic fiber \(\mathbb G_{a,K}\). The
affine \(R\)-algebra \(\mathcal A(G)\) is the subring of
\(K[X]\) consisting of the polynomials whose constant
coefficient lies in \(R\).
\textup{(N.B. We already encountered this example in
\hyperref[rem:sga3-VIB-11.10.1]{Remark 11.10.1}.)}
Furthermore, \(G\) represents the functor assigning to every
\(R\)-algebra \(B\) the set of sequences
\((x_n)_{n\in\mathbb N}\) of elements of \(B\) such that}
\[
  x_n=\pi x_{n+1}\qquad\text{for every \(n\)}.
\]
\emph{\textup{(In particular, every \(x_n\) is indefinitely
\(\pi\)-divisible.)}}

Now let \(S\) be a Noetherian scheme such that every coherent
\(\mathcal O_S\)-module is a quotient of a locally free
\(\mathcal O_S\)-module of finite type. This holds, for
example, when \(S\) is a separated Noetherian regular scheme;
see SGA~6, Expos\'e~II, 2.1.1 and 2.2.7. \textup{(One can
show that every Noetherian regular scheme of dimension at
most \(1\) also has this property. It does not hold, however,
when \(S\) is the affine plane over a field \(k\) with the
origin doubled; see \emph{loc.\ cit.}, 2.2.7.2.)}

\medskip
\phantomsection\label{def:sga3-VIB-13.4}%
\noindent\textbf{Definition 13.4.---}
\emph{Let \(G\) be an affine flat \(S\)-group. Following
R.~W.~Thomason \textup{(\hyperref[bib:sga3-VIB-Th87]{[Th87]},
\SGARefLink{prop:sga3-VIB-2.1}{2.1})}, we say that the pair \((G,S)\) has the
\emph{equivariant resolution property}, or satisfies
\textup{(RE)}, if, for every coherent
\(G\)-\(\mathcal O_S\)-module \(\mathcal F\), there are a
locally free \(G\)-\(\mathcal O_S\)-module \(\mathcal E\) of
finite rank and a \(G\)-equivariant epimorphism}
\[
  \mathcal E\longrightarrow\mathcal F.
\]

In \emph{loc.\ cit.}, \SGARefLink{def:sga3-VIB-3.1}{Theorem~3.1}, Thomason proves the result
below under the hypothesis that \(G\) be essentially free
over \(S\) \textup{(see the remark below)}. Gabber pointed
out to us the simpler proof given here, which does not use
that hypothesis.

\medskip
\phantomsection\label{prop:sga3-VIB-13.5}%
\noindent\textbf{Proposition 13.5.---}
\emph{Let \(S\) be a Noetherian scheme, and let \(G\) be an
affine \(S\)-group, flat and of finite type, with affine
algebra \(\mathcal A(G)\). Suppose that \((G,S)\) satisfies
\textup{(RE)}. Then:}
\begin{enumerate}
\item[\textup{(i)}]\phantomsection\label{concl:sga3-VIB-13.5-i}
  \emph{\(G\) is isomorphic to a closed subgroup of}
  \[
    H=\operatorname{Aut}_{\mathcal O_S}(\mathcal V)
  \]
  \emph{for some locally free \(\mathcal O_S\)-module
  \(\mathcal V\) of finite rank.}

\item[\textup{(ii)}]\phantomsection\label{concl:sga3-VIB-13.5-ii}
  \emph{If, moreover, \(S\) is affine, one may take
  \(\mathcal V=\mathcal O_S^{\oplus d}\) for some \(d\), so
  that \(H=\operatorname{GL}_{d,S}\).}%
  \begingroup
  \renewcommand{\thefootnote}{\fnsymbol{footnote}}%
  \footnote[2]{\label{note:sga3-VIB-13.5-ii-rank-letter}
  Translator's note: the French statement prints ``for some
  \(n\)'' after \(\mathcal O_S^{\oplus d}\), while its conclusion
  is \(H=\operatorname{GL}_{d,S}\). The same rank is denoted by
  \(d\) here. Reinhold's English translation likewise uses
  \(d\).}%
  \endgroup
\end{enumerate}

The proof is analogous to that of
\hyperref[prop:sga3-VIB-13.2]{13.2}. As there, there is a
coherent \(G\)-stable \(\mathcal O_S\)-submodule
\(\mathcal F\) that generates
\(\mathcal A=\mathcal A(G)\) as an
\(\mathcal O_S\)-algebra. Replacing \(S\) by one of its
connected components, we may suppose that \(S\) is
connected. By \textup{(RE)}, there are a locally free
\(G\)-\(\mathcal O_S\)-module \(\mathcal E\) of rank \(n\)
and an epimorphism of \(\mathcal A\)-comodules
\[
  \pi:\mathcal E\longrightarrow\mathcal F.
\]
Put \(H=\operatorname{Aut}_{\mathcal O_S}(\mathcal E)\).
This is an \(S\)-group scheme locally isomorphic to
\(\operatorname{GL}_{n,S}\). The action of \(G\) on
\(\mathcal E\) induces a morphism of affine \(S\)-groups
\[
  \rho:G\longrightarrow H,
\]
corresponding to a morphism of \(\mathcal O_S\)-Hopf algebras
\[
  \varphi:\mathcal A(H)\longrightarrow\mathcal A(G).
\]
We shall show that \(\rho\) is a closed immersion.

The question is local on \(S\), so suppose that
\(S=\operatorname{Spec}(R)\) and that
\[
  V=\Gamma(S,\mathcal E)
\]
is a free \(R\)-module with basis \(v_1,\ldots,v_n\). Then
\(H\simeq\operatorname{GL}_{n,R}\), and
\[
  B=\Gamma(S,\mathcal A(H))
\]
is generated as an \(R\)-algebra by the matrix coefficients
\(c_{ij}\) and by \(d^{-1}\), where \(d\) denotes the
determinant. Let \(\Delta\)
\textup{(respectively, \(\varepsilon\))} be the
comultiplication \textup{(respectively, augmentation)} of
\[
  A=\Gamma(S,\mathcal A(G)),
\]
let
\[
  \mu:V\longrightarrow V\otimes_R A
\]
be the \(A\)-comodule structure on \(V\), and put
\[
  F=\Gamma(S,\mathcal F).
\]
For \(j=1,\ldots,n\), write
\phantomsection\label{form:sga3-VIB-13.5-mu-matrix}%
\[
  \mu(v_j)=\sum_{i=1}^{n}v_i\otimes a_{ij}.
\]
Then \(\varphi(c_{ij})=a_{ij}\). Since
\(\pi:V\to F\) is a morphism of \(A\)-comodules, one has
\phantomsection\label{form:sga3-VIB-13.5-comodule-map}%
\[
  \sum_{i=1}^{n}\pi(v_i)\otimes a_{ij}
    =
  (\pi\otimes_R\mathrm{id}_A)(\mu(v_j))
    =
  \Delta(\pi(v_j)),
\]
and therefore
\phantomsection\label{form:sga3-VIB-13.5-coefficient-recovery}%
\[
  \begin{aligned}
  \pi(v_j)
    &=(\varepsilon\otimes_R\mathrm{id}_A)\Delta(\pi(v_j))\\
    &=\sum_{i=1}^{n}\varepsilon(\pi(v_i))a_{ij}\\
    &=\varphi\left(
        \sum_{i=1}^{n}\varepsilon(\pi(v_i))c_{ij}
      \right).
  \end{aligned}
\]
Thus \(\varphi(B)\) contains \(\pi(V)=F\), which generates
\(A\) as an \(R\)-algebra. Hence \(\varphi\) is surjective,
and \(\rho\) is a closed immersion. This proves
\hyperref[concl:sga3-VIB-13.5-i]{\textup{(i)}}.

If \(S\) is affine, there is a locally free
\(\mathcal O_S\)-module \(\mathcal V'\) of finite rank such
that
\[
  \mathcal V\oplus\mathcal V'
    =\mathcal O_S^{\oplus d}.
\]
Regard \(\mathcal V'\) as a trivial \(G\)-module and replace
\(\mathcal V\) by \(\mathcal O_S^{\oplus d}\). This realizes
\(G\) as a closed subgroup of \(\operatorname{GL}_{d,S}\)
and completes the proof.

\medskip
\phantomsection\label{rem:sga3-VIB-13.6}%
\noindent\textbf{Remarks 13.6.---}
\begin{enumerate}
\item[\textup{(a)}]\phantomsection\label{rem:sga3-VIB-13.6-a}
  For completeness, let us briefly sketch Thomason's
  argument \textup{(\hyperref[bib:sga3-VIB-Th87]{[Th87]},
  \SGARefLink{def:sga3-VIB-3.1}{Theorem~3.1})}, retaining the preceding notation. The
  \(G\)-equivariant epimorphism
  \[
    \pi:\mathcal E\longrightarrow\mathcal F
  \]
  induces a closed immersion
  \[
    \tau:G\longrightarrow V(\mathcal E)
  \]
  such that
  \[
    \tau(g'g)=\tau(g')\cdot g.
  \]
  \textup{(N.B. \(G\) acts on the right on
  \(V(\mathcal E)\).)}
  There is an isomorphism
  \[
    H\simeq
    \operatorname{Aut}_{\mathcal O_S}
      \bigl(V(\mathcal E)\bigr)^{\mathrm{op}}
  \]
  compatible with the left action of \(G\) on
  \(\mathcal E\) and its right action on \(V(\mathcal E)\).
  Let
  \[
    N=\operatorname{Transp}^{\mathrm{str}}_H
      \bigl(\tau(G),\tau(G)\bigr)
  \]
  be the strict transporter. If \(G\) is essentially free
  over \(S\), it follows from
  \hyperref[ex:sga3-VIB-6.2.4-e]{6.2.4(e)} that \(N\) is a
  closed subgroup scheme of \(H\), hence affine over \(S\).
  Moreover, \(\rho\) factors through a morphism of
  \(S\)-groups
  \[
    \rho':G\longrightarrow N.
  \]
  For every \(S'\to S\) and \(h\in N(S')\), put
  \[
    \pi(h)=\tau(1)\cdot h,
  \]
  where \(1\) is the unit element of \(G(S')\). This defines
  a morphism of \(S\)-schemes
  \[
    \pi:N\longrightarrow\tau(G)
  \]
  that is a retraction of \(\rho'\), after identifying
  \(G\) with \(\tau(G)\). Since \(N\) is separated over
  \(S\), it follows that \(\rho'\) is a closed immersion.

\item[\textup{(b)}]\phantomsection\label{rem:sga3-VIB-13.6-b}
  The proof of \hyperref[bib:sga3-VIB-Th87]{[Th87]},
  \SGARefLink{def:sga3-VIB-3.1}{Theorem~3.1}, appears to require the hypothesis that \(G\)
  be essentially free over \(S\), although that hypothesis
  is not stated there \textup{(the author invokes instead
  the fact that \(H\) is essentially free)}. It is,
  however, satisfied when \(G\) is reductive
  \textup{(cf.\ \SGARefLink{sga3:xxii:cor:5.7.8}{Expos\'e~XXII, 5.7.8})}, and hence in all
  cases considered in \emph{loc.\ cit.}, \SGARefLink{rem:sga3-VIB-3.2}{Corollary~3.2}.
  In particular, Thomason proves in \emph{loc.\ cit.}, \SGARefLink{prop:sga3-VIB-2.5}{2.5},
  that if \(S\) is separated, Noetherian, and regular of
  dimension at most \(2\), if \(G\) is affine and of finite
  presentation, and if \(\mathcal A(G)\) is a locally
  projective \(\mathcal O_S\)-module, then \((G,S)\)
  satisfies \textup{(RE)}. By
  \hyperref[prop:sga3-VIB-13.5]{13.5}, this gives
  \hyperref[prop:sga3-VIB-13.2]{13.2} under a slightly more
  restrictive hypothesis.
\end{enumerate}

To conclude, let us point out that the proof of
\hyperref[bib:sga3-VIB-Th87]{[Th87]}, \SGARefLink{prop:sga3-VIB-2.5}{2.5}, can be simplified
slightly as follows. \textup{(For brevity, we work in the
affine situation \(S=\operatorname{Spec}(R)\).)}

\medskip
\phantomsection\label{prop:sga3-VIB-13.7}%
\noindent\textbf{Proposition 13.7.---}
\emph{Let \(R\) be a Noetherian regular ring of dimension at
most \(2\), let \(C\) be an \(R\)-coalgebra that is
projective as an \(R\)-module, and let \(F\) be a
\(C\)-comodule of finite type over \(R\). Then \(F\) is a
quotient of a \(C\)-comodule \(V\) that is projective and of
finite type over \(R\).}

\emph{Proof.}
Replacing \(\operatorname{Spec}(R)\) by one of its connected
components, we may suppose that \(R\) is integral, with
fraction field \(K\). Let \(\Delta\)
\textup{(respectively, \(\varepsilon\))} denote the
comultiplication \textup{(respectively, augmentation)} of
\(C\), and let
\[
  \rho:F\longrightarrow F\otimes C
\]
be the comodule structure on \(F\). Let
\[
  \pi:W\longrightarrow F
\]
be a surjective morphism, with \(W\) a finite free
\(R\)-module. Give \(W\otimes C\) the comodule structure
defined by \(\mathrm{id}_W\otimes\Delta\), and do the same for
\(F\otimes C\). Then
\(\rho:F\to F\otimes C\) is a morphism of \(C\)-comodules
with retraction \(\mathrm{id}_F\otimes\varepsilon\).

Let \(W'\) be the \(C\)-comodule defined by the following
cartesian square:
\begin{center}
  \phantomsection\label{diag:sga3-VIB-13.7-cartesian-square}%
  \begin{tikzcd}[row sep=large,column sep=large]
    W' \arrow[r] \arrow[d] &
    F \arrow[d,"\rho"] \\
    W\otimes C \arrow[r,"\pi\otimes\mathrm{id}_C"] &
    F\otimes C
  \end{tikzcd}
\end{center}
Thus \(W'\) identifies with the kernel of the morphism
\[
  W\otimes C\longrightarrow(F\otimes C)/\rho(F),
\]
and the projection
\[
  \pi':W'\longrightarrow F,
  \qquad
  x\longmapsto(\pi\otimes\varepsilon)(x),
\]
is surjective. Since \(F\) is a finitely generated
\(R\)-module, there is a \(C\)-subcomodule \(V'\) of \(W'\),
finitely generated over \(R\), such that
\[
  \pi'(V')=F.
\]
Since \(W\otimes C\) is \(R\)-torsion-free, so is \(V'\).
Replacing \(F\) by \(V'\), we may therefore suppose from
the outset that \(F\) is torsion-free.

Apply the preceding construction to this new \(F\), obtaining
\(V'\) as above. Consider the subcomodule \(V\) that is the
kernel of the morphism
\phantomsection\label{form:sga3-VIB-13.7-E-quotient}%
\[
  W\otimes C\longrightarrow
  E=\frac{W\otimes C\otimes K}{V'\otimes K}.
\]
Then \(V\) contains \(V'\), and
\[
  Q=(W\otimes C)/V
\]
is an \(R\)-submodule of the \(K\)-vector space \(E\). Put
\[
  Q'=E/Q.
\]
Since \(W\otimes C\) and \(E\) are flat \(R\)-modules, for
every \(R\)-module \(N\) one obtains
\phantomsection\label{form:sga3-VIB-13.7-Tor}%
\[
  \operatorname{Tor}_1^R(V,N)
    \simeq
  \operatorname{Tor}_2^R(Q,N)
    \simeq
  \operatorname{Tor}_3^R(Q',N).
\]
Because \(R\) is regular of dimension at most \(2\), the
right-hand term vanishes. Hence \(V\) is a flat
\(R\)-module.

It remains to prove that \(V\) is finitely generated. Put
\[
  M=W\otimes C.
\]
By construction,
\[
  V/V'\simeq(M/V')_{\mathrm{tors}},
\]
the \(R\)-torsion submodule of \(M/V'\). Since \(M\) is a
projective \(R\)-module, there is a projective \(R\)-module
\(P\) such that
\[
  M\oplus P=L
\]
is free. Then
\[
  (M/V')_{\mathrm{tors}}
    \simeq
  (L/V')_{\mathrm{tors}}.
\]
Since \(V'\) is finitely generated, there is a direct
summand
\[
  L'\simeq R^n
\]
of \(L\) containing \(V'\). Hence
\[
  (L/V')_{\mathrm{tors}}
    \simeq
  (L'/V')_{\mathrm{tors}},
\]
and the latter is finitely generated because \(L'/V'\) is.
Thus \(V\) is a finitely generated flat \(R\)-module, hence
a finitely generated projective \(R\)-module, since \(R\) is
Noetherian. This proves \hyperref[prop:sga3-VIB-13.7]{Proposition~13.7}.

\section*{Bibliography}
\phantomsection\label{bib:sga3-VIB}%
\addcontentsline{toc}{section}{Bibliography}
\footnote{\label{note:sga3-VIB-140}Editor's note: additional
references cited in this Expos\'e.}

\begin{description}
\item[{[An73]}]\phantomsection\label{bib:sga3-VIB-An73}%
  S.~Anantharaman,
  \emph{Sch\'emas en groupes, espaces homog\`enes et espaces
  alg\'ebriques sur une base de dimension 1},
  M\'em. Soc. Math. France \textbf{33} (1973), 5--79.
\item[{[Ba55]}]\phantomsection\label{bib:sga3-VIB-Ba55}%
  I.~Barsotti,
  \emph{Un teorema di struttura per le variet\`a gruppali},
  Atti Acad. Naz. Lincei Rend. Cl. Sci. Fis. Mat. Nat.
  \textbf{18} (1955), 43--50.
\item[{[BLR]}]\phantomsection\label{bib:sga3-VIB-BLR}%
  S.~Bosch, W.~L\"utkebohmert, and M.~Raynaud,
  \emph{N\'eron Models}, Springer-Verlag, 1990.
\item[{[BAC]}]\phantomsection\label{bib:sga3-VIB-BAC}%
  N.~Bourbaki,
  \emph{Alg\`ebre commutative}, Chapters I--IV, V--VII,
  and X, Masson, 1985 and 1998.
\item[{[BTop]}]\phantomsection\label{bib:sga3-VIB-BTop}%
  N.~Bourbaki,
  \emph{Topologie g\'en\'erale}, Chapters I--IV, Hermann,
  1971.
\item[{[Br09]}]\phantomsection\label{bib:sga3-VIB-Br09}%
  M.~Brion,
  \emph{Anti-affine algebraic groups},
  J. Algebra \textbf{321} (2009), no.~3, 934--952.
\item[{[Ch60]}]\phantomsection\label{bib:sga3-VIB-Ch60}%
  C.~Chevalley,
  \emph{Une d\'emonstration d'un th\'eor\`eme sur les groupes
  alg\'ebriques},
  J. Maths. Pure Appl. \textbf{39} (1960), 307--317.
\item[{[Co02]}]\phantomsection\label{bib:sga3-VIB-Co02}%
  B.~Conrad,
  \emph{A modern proof of Chevalley's theorem on algebraic
  groups},
  J. Ramanujan Math. Soc. \textbf{17} (2002), 1--18.
\item[{[DG70]}]\phantomsection\label{bib:sga3-VIB-DG70}%
  M.~Demazure and P.~Gabriel,
  \emph{Groupes alg\'ebriques}, Masson \& North-Holland,
  1970.
\item[{[Gr73]}]\phantomsection\label{bib:sga3-VIB-Gr73}%
  L.~Gruson,
  \emph{Dimension homologique des modules plats sur un anneau
  noeth\'erien},
  Symposia Mathematica \textbf{XI} (1973), 243--254.
\item[{[Kn71]}]\phantomsection\label{bib:sga3-VIB-Kn71}%
  D.~Knutson,
  \emph{Algebraic spaces},
  Lecture Notes in Mathematics \textbf{203},
  Springer-Verlag, 1971.
\item[{[Ma07]}]\phantomsection\label{bib:sga3-VIB-Ma07}%
  B.~Margaux,
  \emph{Passage to the Limit in Non-Abelian \v{C}ech
  Cohomology},
  J. Lie Theory \textbf{17} (2007), 591--596.
\item[{[MW94]}]\phantomsection\label{bib:sga3-VIB-MW94}%
  A.~Masuoka and D.~Wigner,
  \emph{Faithful flatness of Hopf algebras},
  J. Algebra \textbf{170} (1994), 156--184.
\item[{[Per76]}]\phantomsection\label{bib:sga3-VIB-Per76}%
  D.~Perrin,
  \emph{Approximation des sch\'emas en groupes,
  quasi-compacts sur un corps},
  Bull. Soc. Math. France \textbf{104} (1976), 323--335.
\item[{[Pes66]}]\phantomsection\label{bib:sga3-VIB-Pes66}%
  C.~Peskine,
  \emph{Une g\'en\'eralisation du ``Main Theorem'' de Zariski},
  Bull. Sci. Math. \textbf{90} (1966), 119--127.
\item[{[PY06]}]\phantomsection\label{bib:sga3-VIB-PY06}%
  G.~Prasad and J.-K.~Yu,
  \emph{On quasi-reductive group schemes},
  J. Algebraic Geom. \textbf{15} (2006), 507--549.
\item[{[Ray70a]}]\phantomsection\label{bib:sga3-VIB-Ray70a}%
  M.~Raynaud,
  \emph{Faisceaux amples sur les sch\'emas en groupes et les
  espaces homog\`enes},
  Lecture Notes in Mathematics \textbf{119},
  Springer-Verlag, 1970.
\item[{[Ray70b]}]\phantomsection\label{bib:sga3-VIB-Ray70b}%
  M.~Raynaud,
  \emph{Anneaux locaux hens\'eliens},
  Lecture Notes in Mathematics \textbf{169},
  Springer-Verlag, 1970.
\item[{[RG71]}]\phantomsection\label{bib:sga3-VIB-RG71}%
  M.~Raynaud and L.~Gruson,
  \emph{Crit\`eres de platitude et de projectivit\'e},
  Invent. Math. \textbf{13} (1971), 1--89.
\item[{[Ro56]}]\phantomsection\label{bib:sga3-VIB-Ro56}%
  M.~Rosenlicht,
  \emph{Some basic theorems on algebraic groups},
  Amer. J. Math. \textbf{78} (1956), 401--443.
\item[{[SS09]}]\phantomsection\label{bib:sga3-VIB-SS09}%
  C.~Sancho de Salas and F.~Sancho de Salas,
  \emph{Principal bundles, quasi-abelian varieties and
  structure of algebraic groups},
  J. Algebra \textbf{322} (2009), no.~8, 2751--2772.
\item[{[Se68]}]\phantomsection\label{bib:sga3-VIB-Se68}%
  J.-P.~Serre,
  \emph{Groupes de Grothendieck des sch\'emas en groupes
  r\'eductifs d\'eploy\'es},
  Publ. Math. I.H.\'E.S. \textbf{34} (1968), 37--52.
\item[{[Se99]}]\phantomsection\label{bib:sga3-VIB-Se99}%
  C.~S.~Seshadri,
  \emph{Chevalley: some reminiscences},
  Transform. Groups \textbf{4} (1999), nos.~2--3, 119--125.
\item[{[Ta72]}]\phantomsection\label{bib:sga3-VIB-Ta72}%
  M.~Takeuchi,
  \emph{A correspondence between Hopf ideals and sub-Hopf
  algebras},
  Manuscripta Math. \textbf{7} (1972), 251--270.
\item[{[Ta79]}]\phantomsection\label{bib:sga3-VIB-Ta79}%
  M.~Takeuchi,
  \emph{Relative Hopf modules---Equivalences and freeness
  criteria},
  J. Algebra \textbf{60} (1979), 452--471.
\item[{[Th87]}]\phantomsection\label{bib:sga3-VIB-Th87}%
  R.~W.~Thomason,
  \emph{Equivariant resolution, linearization, and Hilbert's
  fourteenth problem over arbitrary base schemes},
  Adv. Math. \textbf{65} (1987), 16--34.
\end{description}

